\begin{example}
    Sean $(x_n)_{n\geq1},(y_n)_{n\geq1}$ dos sucesiones equidistribuidas módulo $1$.
    Entonces la sucesión
    \[
        z_n = \begin{cases}
            x_k & \text{si } n = 2k-1,\\
            y_k & \text{si } n = 2k,
        \end{cases}
    \]
    también es equidistribuida módulo $1$.
\end{example}

\begin{resolve}
    Sea un intervalo $(a,b) \subset \mathbb{T}$. Consideremos $2N$ términos de la sucesión $(z_n)$, 
    \[
        A_{2N}((a,b), z) = A_N((a,b), x) + A_N((a,b), y)
    \]
    
    Teniendo en cuenta que $(x_n)$ e $(y_n)$ son equidistribuidas módulo $1$, al dividir por $2N$
    y pasar al límite, obtenemos
    \[
        \lim_{N \to \infty} \frac{A_{2N}((a,b), z)}{2N} =
        \lim_{N \to \infty} \frac{1}{2} \left( \frac{A_N((a,b), x)}{N} + \frac{A_N((a,b), y)}{N} \right) =
        b-a
    \]
    
    Por otra parte, si tomamos $2N+1$ términos de la sucesión $(z_n)$, 
    \[
        A_{2N+1}((a,b), z) = A_{2N}((a,b), z) + \mathbf{1}_{(a,b)}(x_{N+1})
    \]

    Como $0 \leq \mathbf{1}_{(a,b)}(x_{N+1}) \leq 1$, al dividir por $2N+1$ y tomando límites,
    \[
        \lim_{N \to \infty} \frac{A_{2N+1}((a,b), z)}{2N+1} =
        \lim_{N \to \infty} \left(\frac{2N}{2N+1} \frac{A_{2N}((a,b), z)}{2N} + \frac{\mathbf{1}_{(a,b)}(x_{N+1})}{2N+1} \right) =  
        b-a
    \]

    Por lo tanto, $(z_n)$ es equidistribuida módulo $1$.  
\end{resolve}
